
%\documentclass[oneside,final,14pt]{extreport}
\documentclass[12pt, a4paper, openany]{book}
\usepackage[left=2cm,right=2cm,top=1.5cm,bottom=1.5cm,bindingoffset=0cm]{geometry}
%\usepackage[koi8-r]{inputenc}
\usepackage[russianb]{babel}
\usepackage{vmargin}
\setpapersize{A4}
\usepackage[T2A]{fontenc}
\usepackage[utf8x]{inputenc} % more recent versions (at least>=2004-17-10)
%\usepackage[russian]{babel}
%\setmarginsrb{2cm}{1.5cm}{1cm}{1.5cm}{0pt}{0mm}{0pt}{13mm}
\usepackage{indentfirst}
\usepackage{nicefrac} % For comparison
%\usepackage{xfrac} % Works better with other fonts
%\usepackage[unicode, pdftex]{hyperref}
\usepackage{lettrine}
\usepackage[usenames]{color}
%\special{papersize=a5}
\usepackage{colortbl}
\usepackage[pagestyles]{titlesec}
\usepackage{xfrac} % Works better with other fonts
\usepackage[colorlinks=true,linkcolor=black,urlcolor=black,bookmarksopen=true]{hyperref}


% Настройка вертикальных и горизонтальных отступов
\titlespacing{\chapter}{0pt}{5pt}{5pt}
\titlespacing{\section}{\parindent}{4mm}{4mm}
\titlespacing{\subsection}{\parindent}{3mm}{3mm}

\newcommand{\anonsection}[1]{ \section*{#1} \addcontentsline{toc}{section}{\numberline {}#1}} 

\makeatletter %%%%% <---- Starting chapter without a pagebreak
\renewcommand\chapter{\par%
	\thispagestyle{plain}% \global\@topnum\z@
	\@afterindentfalse \secdef\@chapter\@schapter}
\makeatother %%%%% <---- Starting chapter without a pagebreak
\titleformat{\chapter}[display]
{\normalfont\bfseries}{}{0pt}{\Large}

\newpagestyle{mystyle}{
	\sethead[\thepage][][]{}{}{\thepage}	
}

\pagestyle{mystyle}

\sloppy
\begin{document}

	
	\begin{titlepage}
		
		\begin{center}
			\topskip0pt
			\vspace*{\fill}
			
			
			{\large\bf Ковалев В. А.\\}
			\ \\
			\ \\
			{\Huge\bf КРУПНЕЙШИЕ УГОЛОВНЫЕ ДЕЛА ХХ ВЕКА В США\\}
			\ \\
			\ \\
			\ \\
			\vspace*{\fill}
			
			\vfill
			
			
			МОСКВА
			
			«ЮРИДИЧЕСКАЯ ЛИТЕРАТУРА»
			
			1990
			

		\end{center}
		
	\end{titlepage}
	
	\thispagestyle{empty} % выключаем отображение номера для этой страницы
	
	\newpage
	

\setcounter{secnumdepth}{0}
	
	\section[Вместо предисловия]{\center \textbf{Вместо предисловия}}
	
Эта книга не научный трактат по юриспруденции и не детективный роман с пикантными подробностями леденящих душу преступлений. В ней рассказывается о  подлинных уголовных делах, которые потрясли Америку XX века. Вокруг судебных процессов по этим уголовным делам разгорелись не только юридические дискуссии, но  бушевали и политические страсти. Наряду с профессиональными юристами в них принимали участие видные  государственные деятели, сенаторы, министры, президенты США{\footnote{1}}. Едва ли можно было найти американца, который просмотр утренних газет не начинал бы с поисков новой информации о ходе расследования и судебного разбирательства этих дел. Люди выходили на массовые митинги и манифестации; в адрес губернаторов штатов и непосредственно в Белый дом направлялись многочисленные требования и протесты. Ни один американец не оставался равнодушным к судьбам лиц, сидящих на скамье подсудимых. 

Чем же был вызван такой пристальный интерес к этим судебным процессам? Ведь американцев не удивишь  сообщениями о таинственных убийствах и групповых  изнасилованиях, вооруженных нападениях и похищениях людей, 
организации взрывов и головоломных хищениях  
астрономических сумм. И даже судебные процессы об  
антиправительственных заговорах и атомном шпионаже в пользу 
СССР — не столь уж редкие события в послевоенной 
истории американской юстиции. Все это нашло отражение 
и в уголовных делах, собранных в книге, которую мы 
предлагаем читателям. Однако неослабевающий интерес 
к этим судебным процессам вызван отнюдь не сюжетной 
занимательностью дел, а обусловлен причинами,  
несравненно более важными и значительными. Главная из них 
состоит в том, что приговор по этим уголовным делам, 
1» 
Крупнейшие уголовные дела XX века в США 
4 
строго говоря, касался не только подсудимых. В условиях 
прецедентной системы, когда каждое решение высших 
судебных инстанций приобретает нормативное значение 
и может быть использовано в качестве источника права 
в последующих процессах, судебные прения между  
обвинением и защитой по уголовному делу неизбежно  
превращаются в арену борьбы по вопросам неотъемлемых прав 
человека и основных гражданских свобод. 
Большое внимание в книге уделяется деятельности 
суда присяжных в США. Эта проблематика в последнее 
время приобрела актуальность и в нашей стране в связи 
с многочисленными предложениями со стороны  
профессиональных юристов и общественности о введении суда 
присяжных в СССР. Поэтому зарубежный опыт (как 
положительный, так и отрицательный) представляет  
существенный интерес. Читатель на материалах крупнейших 
уголовных дел в США познакомится с практикой подбора 
присяжных и их отвода, узнает о борьбе мнений среди 
членов жюри и давлении на них со стороны власть 
предержащих. И хотя тайна совещания судей охраняется 
законом и разглашению не подлежит, тем не менее 
некоторые приведенные в книге данные из американских 
источников позволят составить достаточно ясное  
представление о том, что происходит в совещательной комнате, 
которую присяжные заседатели не вправе покинуть до 
вынесения вердикта по уголовному делу. 
Читатель получит представление не только о  
торжественной процедуре в помпезных залах дворцов правосудия, 
но и познакомится с закулисной стороной производства по 
уголовным делам, которая обычно остается вне поля 
зрения непосвященных. 
В самом центре американской столицы — как раз 
между Белым домом и Капитолием — расположилось 
новое грандиозное здание Федерального бюро  
расследований. Более полувека это ведомство возглавлял Эдгар Гувер, 
известный также по кличке Револьвер № 1, которой 
окрестили его американцы за патологическую страсть 
к преследованиям инакомыслящих, привычку сразу же 
хвататься за оружие при появлении «красной опасности». 
И после смерти Гувера его зловещий лозунг «инакомыслие 
равносильно предательству» не утратил своего значения 
в деятельности американских спецслужб. 
Сегодня в Соединенных Штатах нередко вспоминают 
Вместо предисловия 
5 
старину Эйба (Авраама Линкольна. — В. К.)у который 
«смотрел в корень, когда говорил, что в Америке все 
считают себя доверенными лицами свободы, но при 
употреблении этого слова вкладывают в него разный 
смысл» 2. В принципе, конечно, «каждый американец может 
выйти на улицу, встать на ящик из-под мыла и высказать 
людям все, что у него наболело, — пишет известный 
американский историк, профессор Говард Зинн. — Но надо 
долго ждать, когда Конституция тебя защитит, прямую же 
власть над твоим «свободным» словом имеют те, у кого 
в руках дубинка и револьвер в кобуре и кто обязательно 
появится в том самом месте, где ты этой «свободой» 
захочешь воспользоваться» \ 
Американский профессор истории знает, о чем говорит. 
Ведь наука, которую он представляет, убедительно  
свидетельствует, что многие страницы современной истории 
Соединенных Штатов отмечены и полицейскими  
расправами над «радикалами», и судебными преследованиями 
инакомыслящих. 
В уголовных процессах против этой категории  
американских граждан сконцентрировались классовые  
противоречия и социальные столкновения эпохи. Собранные в этой 
книге судебные дела — вехи политической истории США 
XX века. Когда вопреки Конституции страны уголовному 
преследованию за политические убеждения подвергается 
гражданин Соединенных Штатов — страны, претендующей 
на роль лидера «свободного общества», — причины этого, 
если взглянуть на поверхность данного явления, можно 
отнести на счет злоупотреблений спецслужб,  
безответственности судей, предвзятости специально подобранных 
присяжных. Однако рассмотренные в совокупности — 
в историческом контексте и с учетом борьбы  
противоположных интересов — эти судебные процессы позволяют 
увидеть глубинную политическую природу уголовных дел 
против демократических сил страны — борцов за  
гражданские права, активистов антивоенного движения,  
профсоюзных лидеров. За разнообразными частными юридическими 
особенностями каждого дела отчетливо просматривается 
единая политическая линия на ограничение социальной 
активности инакомыслящих и инакодействующих. 
В настоящей книге собрана, разумеется, лишь  
небольшая часть крупных судебных процессов, которые имеют 
политические основания. Автор предлагаемой читателям 
Крупнейшие уголовные дела XX века в США 
6 
книги не ставил перед собой задачу исследовать историю 
американской уголовной юстиции XX века. Его намерения 
более скромны — написать серию очерков, которые бы 
давали наглядное представление о крупных судебных 
делах, проходивших по ведомству уголовной юстиции, но 
имевших очевидную политическую подоплеку. 
Многие из рассмотренных на страницах книги  
судебных процессов представляют значительный юридический 
интерес с точки зрения уголовного и  
уголовно-процессуального права. Они неоднократно пересматривались в  
различных апелляционных инстанциях страны, дошли до  
Верховного суда США, породили судебные прецеденты, имеющие 
в условиях англосаксонской системы права нормативное 
значение. Другие приведенные в книге уголовные дела 
с юридической стороны, быть может, более очевидны, но 
интересны фактической фабулой, системой доказательств, 
широким общественным резонансом, который они вызвали 
в США и за рубежом. За некоторыми из них внимательно 
следил Владимир Ильич Ленин, что нашло отражение в его 
трудах. 
В книгу вошли и относительно скромные по своей 
фабуле и юридическим особенностям судебные дела, 
которые некогда прошли незамеченными, но, как оказалось 
впоследствии, имели колоссальное влияние на дальнейшее 
развитие прецедентного права в США. Возникла  
необходимость рассмотреть значительное число уголовных дел, 
ранее не освещавшихся на страницах советской  
юридической литературы и в силу этого малоизвестных или совсем 
неизвестных нашему читателю, потребовавшая  
использования ряда новых документальных материалов, которые на 
русском языке публикуются впервые. С целью сохранения 
лексики и стилистики подлинников диалоги из протоколов 
судебных заседаний приводятся дословно. В них по 
возможности сохранены не только фактическая сторона, 
но и эмоциональная окраска показаний, заявлений,  
выступлений в прениях и других высказываний участников 
процесса. 
Отдельные судебные процессы описаны скупо, подчас 
даже лапидарно, и связано это с тем, что некоторые 
документальные материалы не сохранились, другие  
хранятся в архивах спецслужб и едва ли в обозримом будущем 
станут достоянием гласности. Поэтому освещая лишь 
твердо установленные факты, будем сомневаться в других. 
Вместо предисловия 
7 
И еще одно замечание принципиального свойства. 
При оценке юридических особенностей и социального 
значения рассматриваемых уголовных дел автор данной 
книги старался воздерживаться от патетической риторики, 
которая некоторое время преобладала в массовых  
публикациях о буржуазном государстве и праве и которой он сам 
не избежал в своих более ранних публикациях. Из этой 
книги читатель узнает не только о необоснованных  
вердиктах присяжных, незаконных приговорах и неоправданно 
жестоких мерах наказания. Приводятся данные и другого 
характера — об отмене незаконных судебных решений 
первой инстанции Верховными судами штатов или  
Верховным судом США; об оправдании подсудимых,  
привлеченных к уголовной ответственности по сфабрикованным 
обвинениям; об осуждении американским судом  
лжесвидетелей из числа агентов спецслужб за дачу ложных 
показаний. Словом, для этой книги судебные процессы 
отбирались отнюдь не с целью рассмотреть одни только 
негативные явления в системе уголовной юстиции США. По 
мнению автора, более конструктивная задача — показать 
реальные обстоятельства судебных дел во всей их  
сложности и противоречивости, борьбу демократического и  
реакционного начал в состязательном уголовном процессе. 
Книга написана на строго документальных материалах, 
в ней нет вымышленных имен. У автора не было  
необходимости драматизировать события, поскольку реальные  
факты, о которых повествует книга, отражают подчас такие 
человеческие трагедии, перед которыми меркнет самая 
изощренная фантазия. 


\newpage

	\section[Глава первая. Почему вы стали телезрителями? Какие невидимые связи установились между вашими детьми и телевидением?]{\center \textbf{ГЛАВА 1. ПОЧЕМУ ВЫ СТАЛИ ТЕЛЕЗРИТЕЛЯМИ? КАКИЕ НЕВИДИМЫЕ СВЯЗИ УСТАНОВИЛИСЬ МЕЖДУ ВАШИМИ ДЕТЬМИ И ТЕЛЕВИДЕНИЕМ? }}
	
	
	\newpage
	\tableofcontents
	
	\thispagestyle{empty} % 
	
	\newpage
	
	\setcounter{secnumdepth}{0}
	
	\phantomsection
	
		\section*{Описание}
	
	{\bf Название:} А если в семье телевизор... и дети?
	
{\bf Автор:} Ксенофонтов Валентин Васильевич

{\bf Редактор:} О. Г. Свердлова
	
{\bf Издательство:} М.: <<Знание>>. 1973. - 96 с.
	
		{\bf Аннотация:} Ученые всех стран мира сходятся в одном — обилие телевизионной информации, даже полезной по своему содержанию, без правильного руководства со стороны взрослых воспитывает у детей душевную и нравственную пассивность, приучает довольствоваться малым, не говоря уже о том, что оно отвлекает от чтения, спорта и других занятий. И в то же время телевидение представляет детям весь мир наглядно, в ярких динамичных образных формах. Чего только не увидишь по телевизору! Луноход и дно океана, сказку и экзотических животных. 
		
		Книга учит родителей правильно использовать огромные возможности телевидения для воспитания детей. Она рассчитана на самые широкие круги читателей. 
		


		\thispagestyle{empty} % выключаем отображение номера для этой страницы

	
\end{document}


